\documentclass[10pt,a4paper]{amsart}
\usepackage[margin=19mm]{geometry}
\usepackage{amsmath,amssymb,mathtools}
\usepackage[scr=rsfs,cal=euler]{mathalfa}
\usepackage{xfrac}
\usepackage[unicode,hidelinks,bookmarks=false]{hyperref}
\newcommand{\ie}{i.e.\ }
\newcommand{\cf}{cf.\ }
\newcommand{\resp}{respectively\ }
\newcommand{\Subsection}{\subsection}
\providecommand{\nobreakdash}{\nobreak\hbox{-}}
\setlength{\emergencystretch}{2em}
\multlinegap0pt
\usepackage{amssymb}
\usepackage{mathtools}
\usepackage{tikz}
\usepackage{csquotes}
\usepackage[shortlabels]{enumitem}
\usepackage{xcolor}
\newtheorem{lemma}{Lemma}
\def\R{\mathbb{R}}
\newcommand{\abs}[1]{\left|#1\right|}
\title{Closing the gap: Maz'ya-Shaposhnikova and asymptotics of fractional perimeters}
\author{Elisa Davoli, Alberto Fanizza, Marco Picerni}
\date{Source: arXiv:2605.03955v1 (CC BY 4.0)}
\begin{document}
\maketitle

\noindent\emph{Source and licence.} Closing the gap: Maz'ya-Shaposhnikova and asymptotics of fractional perimeters. Version arXiv:2605.03955v1 (CC BY 4.0), distributed by arXiv under the Creative Commons Attribution 4.0 International licence, http://creativecommons.org/licenses/by/4.0/, as recorded in the arXiv licence metadata for this exact version and shown on the abstract page https://arxiv.org/abs/2605.03955v1. Full licence notice: \texttt{licenses/arxiv-2605.03955-CC-BY-4.0.txt}.\par\smallskip

\noindent\textit{Editorial scope.} Counted unit: the article's lemma lemma:lualphau (source0.tex lines 330--341). The article's complete original proof of it is delivered with it (source0.tex lines 342--411, one \texttt{proof} environment of 70 lines). No mathematics is written, repaired or re-derived: the statement and the proof are the article's own lines, in the article's own notation, and no retained line is substituted or re-flowed.  No further source-local context is needed: every cross-reference of every retained line resolves inside the two delivered spans.  Nothing was omitted from any retained span, so the package carries no omission record.  No retained line contains a citation, so the wrapper carries no bibliography and no bibliography entry.  No retained line contains a figure environment, an image-inclusion command or drawing code, so no image is delivered and the package declares no asset dependency.  Editorial formatting only, in addition: an AMS \texttt{amsart} preamble that restates the article's own declarations and macros used by the retained lines, namely the environments \texttt{lemma} on separate counters, together with the standard packages those lines need.  The retained environments print in this document as Lemma 1 in order of appearance, whereas the article prints the same items with the numbering of its own full document.  The complete native source of the article as distributed in the arXiv e-print for this exact version is bundled unchanged as \texttt{sources/199831/source0.tex}.
\begin{lemma}
    \label{lemma:lualphau}
    Let $u\in L^1_{loc}(\R^d)$. Assume that $\alpha(u^+),\,\alpha(u^-)\in \R$. Then,
    \begin{equation}
        \label{luwelldefined}
        \begin{split}
             L_u&:=\lim_{R\to\infty}\liminf_{s\to 0^+}s\int_{B_R^c(x)}\frac{u(y)}{\abs{x-y}^{d+2s}}\,dy\\
             &=\lim_{R\to\infty}\limsup_{s\to 0^+}s\int_{B_R^c(x)}\frac{u(y)}{\abs{x-y}^{d+2s}}\,dy
        \end{split}
    \end{equation}
    is well defined and does not depend on $x\in\R^d$. Moreover, $L_u=\alpha(u)$.
\end{lemma}

\begin{proof}
    We prove that, if $\alpha(u^+)$ and $\alpha(u^-)$ exist and are finite, then for any $x\in\R^d$ both limits in \ref{luwelldefined} exist and are equal to $\alpha(u)$.\par
    \textit{Step $1$.} Let $0<r<R<+\infty$. For any $x\in \R^d$ we have
    \begin{equation*}
        \limsup_{s\to 0^+} \left|s\int_{B_r^c(x)}\frac{u(y)}{\abs{x-y}^{d+2s}}dy-s\int_{B_R^c(x)}\frac{u(y)}{\abs{x-y}^{d+2s}}dy\right|=0.
    \end{equation*}
    In fact:
    \begin{equation*}
        \begin{split}
            \left|s\int_{B_r^c(x)}\frac{u(y)}{\abs{x-y}^{d+2s}}dy-s\int_{B_R^c(x)}\frac{u(y)}{\abs{x-y}^{d+2s}}dy\right|&\leq s\int_{B_r^c(x)\setminus B_R^c(x)}\frac{\abs{u(y)}}{\abs{x-y}^{d+2s}}\,dy\\
            &\leq \frac{s}{r^{d+2s}}\int_{B_r^c(x)\setminus B_R^c(x)}\abs{u(y)}\,dy,
        \end{split}
    \end{equation*}
    which converges to zero as $s\to0^+$ since $u\in L^1_{loc}(\R^d)$. In particular, choosing $x=0$, by the hypothesis on the existence of $\alpha(u^+)$ and $\alpha(u^-)$, we infer that $\alpha(|u|)$ exists as well, and for every $R>0$:
    \begin{equation}
        \label{Eq:alpha1R}
        \alpha(|u|)=\lim_{s\to0^+}{s}\int_{B_1^c(0)}\frac{|u(y)|}{\abs{y}^{d+2s}}\,dy
        =
        \lim_{s\to0^+}{s}\int_{B_R^c(0)}\frac{|u(y)|}{\abs{y}^{d+2s}}\,dy.
    \end{equation}
    
    \textit{Step 2.} We prove that for every $x\in\R^d$, there holds
    \begin{equation}
    \label{claim:step2}
    \lim_{R\to+\infty}\limsup_{s\to 0^+}   \left|s\int_{B_R^c(x)}\frac{u(y)}{\abs{x-y}^{d+2s}}\,dy-s\int_{B_R^c(0)}\frac{u(y)}{\abs{y}^{d+2s}}\,dy\right|=0.
    \end{equation}
    Indeed, we estimate
     \begin{multline*}
          \left|s\int_{B_R^c(x)}\frac{u(y)}{\abs{x-y}^{d+2s}}\,dy-s\int_{B_R^c(0)}\frac{u(y)}{\abs{y}^{d+2s}}\,dy\right|\\ \leq \left|s\int_{B_R^c(x)}\frac{u(y)}{\abs{x-y}^{d+2s}}\,dy-s\int_{B_R^c(0)}\frac{u(y)}{\abs{x-y}^{d+2s}}\,dy\right|\\
            +\left|s\int_{B_R^c(0)}\frac{u(y)}{\abs{x-y}^{d+2s}}\,dy-s\int_{B_R^c(0)}\frac{u(y)}{\abs{y}^{d+2s}}\,dy\right|
            :=I+II.
    \end{multline*}
    We proceed by estimating each summand separately.\\
    For the first term, we have that
    \begin{equation*}
        \begin{split}
            I&=\left|s\int_{B_R^c(x)}\frac{u(y)}{\abs{x-y}^{d+2s}}\,dy-s\int_{B_R^c(0)}\frac{u(y)}{\abs{y-x}^{d+2s}}\,dy\right|\\
            &  \leq s\int_{B_R(x)\Delta B_R(0)}\frac{\abs{u(y)}}{\abs{x-y}^{d+2s}}dy
        \end{split}
    \end{equation*}
    For $R>0$ big enough (depending on $x$), there exist $0<\delta_1<\delta _2$, (depending on $R$), such that $B_R(x)\Delta B_R(0)\subset B_{\delta_2}(0)\cap B_{\delta_1}^c(x)$. 
    
    Hence,
    \begin{equation*}
       I\leq s\int_{B_{\delta_2}(0)\cap B_{\delta_1}^c(x)}\frac{\abs{u(y)}}{\abs{x-y}^{d+2s}}dy \leq \frac{s}{\delta_1^{d+2s}}\int_{B_{\delta_2(0)}}\abs{u(y)}\,dy,
    \end{equation*}
    where the latter quantity converges to zero as $s\to 0^+$.
    For the second term, we rewrite
    \begin{equation*}
        II=\left|s\int_{B_R^c(0)}u(y)\left(\frac{1}{\abs{y-x}^{d+2s}}-\frac{1}{\abs{y}^{d+2s}}\right)dy\right|.
    \end{equation*}
    Let us fix $y\in \R^d$ and call $h_s(x)=\abs{y-x}^{-(d+2s)}$ for $x\neq y$. A direct computation yields
    \begin{equation*}
        \abs{Dh_s(x)}=\frac{(d+2s)}{\abs{y-x}^{d+2s+1}}.
    \end{equation*}
    For $R>2|x|$, we have that $\abs{y-tx}\geq\abs{y}/2$ for every $y\in B_R^c(0)$ and $t\in [0,1]$. Thus,
    \begin{equation*}
        \abs{\frac{1}{\abs{y-x}^{d+2s}}-\frac{1}{\abs{y}^{d+2s}}}=\abs{h_s(x)-h_s(0)}\leq \frac{(d+2s)\abs{x}2^{d+2s+1}}{\abs{y}^{d+2s+1}.}
    \end{equation*}
    Hence,
    \begin{equation}
    \label{eq:estimateII}
        II\leq \abs{x}(d+2s)\,s\int_{B^c_R(0)}\frac{|u(y)|}{\abs{y}^{d+2s+1}}\,dy\leq \frac{\abs{x}(d+2s)}{R}\,s\int_{B_R^c(0)}\frac{|u(y)|}{\abs{y}^{d+2s}}\,dy,
    \end{equation}
    which, by \eqref{Eq:alpha1R}, entails
    \begin{equation*}
        \lim_{R\to+\infty}\lim_{s\to 0^+}II=0.
    \end{equation*}
    This completes the proof of \eqref{claim:step2}. The statement of the lemma follows then by combining Steps 1 and 2.
\end{proof}

\end{document}
